% Schemi degli esercizi (usati sia nell'esercitazione sia nelle soluzioni)
\definecolor{rame}{HTML}{C46A2B}

\newcommand{\lineaCarico}[3]{% #1 = V, #2 = R_L, #3 = R_C
\begin{circuitikz}[scale=0.8, transform shape, line width=0.6pt]
  \rail{(0,4.4)}{$V = #1$}
  \draw (0,4.4) -- (0,4.2) to[R, l_={$R_L = #2$}, i>^={$I$}] (0,2.2) to[short, *-] (0,2.2)
        to[R, l_={$R_C = #3$}] (0,0.2) -- (0,0) node[sground]{};
  \draw (0,2.2) to[short, -o] (1.4,2.2) node[right]{X};
\end{circuitikz}}

\newcommand{\partitoreI}[3]{% #1 = I, #2 = R1, #3 = R2
\begin{circuitikz}[scale=0.8, transform shape, line width=0.6pt]
  \draw (0,4.4) node[ocirc]{} to[short, i>^={$I = #1$}] (0,3.2);
  \draw (-1.6,3.2) -- (1.6,3.2); \draw (0,3.2) node[circ]{};
  \draw (-1.6,3.2) to[short, i>^={$I_1$}] (-1.6,2.6) to[R, l_={$R_1 = #2$}] (-1.6,0.6) -- (-1.6,0.3) -- (1.6,0.3);
  \draw (1.6,3.2) to[short, i>_={$I_2$}] (1.6,2.6) to[R, l^={$R_2 = #3$}] (1.6,0.6) -- (1.6,0.3);
  \draw (0,0.3) node[circ]{} -- (0,0.1) node[sground]{};
\end{circuitikz}}

\newcommand{\sezionePista}[2]{% #1 = w, #2 = t
\begin{tikzpicture}[line width=0.6pt]
  \fill[rame!70] (0,0) rectangle (3,0.6);
  \draw[rame!60!black, line width=0.9pt] (0,0) rectangle (3,0.6);
  \draw[{Latex[length=2mm]}-{Latex[length=2mm]}] (0,-0.3) -- node[below]{$w = #1$} (3,-0.3);
  \draw[{Latex[length=2mm]}-{Latex[length=2mm]}] (3.3,0) -- node[right]{$t = #2$} (3.3,0.6);
\end{tikzpicture}}
